\section{Language}\label{sec:lang}
In this section, we define the abstract syntax of FGFV and its well-formedness requirements.

\subsection{Notation}

Following, for example, \citet{vytiniotis13} and \citet{coughlin14}, we use an overline above mathematical material to
indicate an appropriately punctuated sequence of zero or more repetitions of the material, with a subscript added
to each metavariable in the material. Thus $\set{\rep{\alpha}}$ is equal to $\set{\alpha_1,\alpha_2,\ldots,\alpha_n}$.
Sometimes we write $\bullet$ as an explicit indication of an empty sequence.
Overlines may be nested, in which case they are expanded outermost first, so
``$\bigobb{\lobindup{\rep{\zeta}}{Q}{\rep{\xi}}}$'' expands first to
``$\bigob{\lobindup{\rep{\zeta_1}}{Q_1}{\rep{\xi_1}},\lobindup{\rep{\zeta_2}}{Q_2}{\rep{\xi_2}},\ldots,\lobindup{\rep{\zeta_n}}{Q_n}{\rep{\xi_n}}}$''
and then to
\newcommand\subfixup[1]{{\setbox0\hbox{$\scriptstyle 2$}\hbox to \wd0{$\scriptstyle #1$\hss}}}
\newcommand\subsubfixup[1]{{\setbox0\hbox{$\scriptscriptstyle 2$}\hbox to \wd0{$\scriptscriptstyle #1$\hss}}}
$$\begin{array}{@{}r@{}l@{}}
{\llbracket\;}&\bd{\zeta_{1\,1},\zeta_{1\,2},\ldots,\zeta_{1\,k_1}}\extends{Q_1}\extends\bd{\xi_{1\,1},\xi_{1\,2},\ldots,\xi_{1\,m_1}}, \\
         &\bd{\zeta_{2\,1},\zeta_{2\,2},\ldots,\zeta_{2\,k_2}}\extends{Q_2}\extends\bd{\xi_{2\,1},\xi_{2\,2},\ldots,\xi_{2\,m_2}}, \\
         &\ldots, \\
         &\bd{\zeta_{\subfixup{n}\,1},\zeta_{\subfixup{n}\,2},\ldots,\zeta_{\subfixup{n}\,k_{\subsubfixup{n}}}}\extends{Q_\subfixup{n}}\extends\bd{\xi_{\subfixup{n}\,1},\xi_{\subfixup{n}\,2},\ldots,\xi_{\subfixup{n}\,m_{\subsubfixup{n}}}}\;\rrbracket
\end{array}$$
(conceptually all on one line, broken here for reasons of space).
The copies of each inner overline expand independently, so the sequences into which they expand may have different lengths; thus
``$\bigobb{\lobindup{\rep{\zeta}}{Q}{\rep{\xi}}}$'' could expand into, e.g.,
``$\llbracket\bd{\zeta_{1\,1},\zeta_{1\,2}}\extends{Q_1}\extends\bd{\xi_{1\,1}},
         \bd{}\extends{Q_2}\extends\bd{\xi_{2\,1},\xi_{2\,2},\xi_{2\,3}}\rrbracket$''.

\newcommand\declsep{\hskip 1.5em}
Following a suggestion of Steele~\citeyearpar[0:59:37]{NEW-OLD-LANGUAGE},
we use an underline below submaterial within overlined material to prevent an 
overline from attaching subscripts to the metavariables in the submaterial;
the effect is to require the underlined material to be the same in each copy.
Thus ``$\rep{\jwfdecl[\same{\Delta}]{\psi}}$'' represents the sequence
``$\jwfdecl[\Delta]{\psi_1}\declsep\jwfdecl[\Delta]{\psi_2}\declsep\cdots\declsep\jwfdecl[\Delta]{\psi_n}$''
for any $n \geq 0$.  Similarly
$\rep{\rep{\jwftype[\same{\same{\Delta'}}]{\xi}}}$ expands to
``$\rep{\jwftype[\same{\Delta'}]{\xi_1}}\declsep\rep{\jwftype[\same{\Delta'}]{\xi_2}}\declsep\cdots\declsep\rep{\jwftype[\same{\Delta'}]{\xi_n}}$''
and then to
$$\begin{array}{@{}l@{\declsep}l@{\declsep}l@{\declsep}l@{}}
\jwftype[\Delta']{\xi_{1\,1}}&\jwftype[\Delta']{\xi_{1\,2}}&\cdots&\jwftype[\Delta']{\xi_{1\,m_1}} \\
\jwftype[\Delta']{\xi_{2\,1}}&\jwftype[\Delta']{\xi_{2\,2}}&\cdots&\jwftype[\Delta']{\xi_{2\,m_2}} \\
\ldots \\
\jwftype[\Delta']{\xi_{\subfixup{n}\,1}}&\jwftype[\Delta']{\xi_{\subfixup{n}\,2}}&\cdots&\jwftype[\Delta']{\xi_{\subfixup{n}\,m_{\subsubfixup{n}}}}
\end{array}$$
The overline notation may explicitly bind index variables by using a range expression.
In such cases, the index variables may appear explicitly in the covered material,
and subscripts are not automatically attached to metavariables.  For example,
$\bigl(\rep{m(e_i)}[3\leq i\leq 6]\bigr)$ means $(m(e_3),m(e_4),m(e_5),m(e_6))$.

Metavariables are expanded \emph{after} overlines.  Multiple occurrences of a metavariable with the same
subscript(s) must have the same expansion, but occurrences of a metavariable with different subscript(s)
may be expanded independently.  Thus ``$\bigobb{\lobindup{\rep{\zeta}}{Q}{\rep{\xi}}}$'' might ultimately expand into, e.g.,
``$\llbracket\bd{\TYP{A},\TYP{B}}\extends{\TYP{U}}\extends\bd{\TYP{List}\ob{\TYP{Boolean}}},
         \bd{}\extends{\TYP{V}}\extends\bd{\Any,\TYP{C},\TYP{Number}}\rrbracket$''.

We use horizontal lines as overlines and underlines, but also in the usual Gentzen notation
for inference rules, with premises above the line and the conclusion below.
To help reduce confusion: (1)~we use a slightly heavier and wider line for inference rules;
(2)~the line for inference rules has more space above and below it;
(3)~we do not use the inference rule notation for axioms (rules that have no premises); and
(4)~we do not nest the inference rule notation to present derivation trees.



\subsection{Syntax}\label{sec:grammar}
\begin{figure}[t]
\leavevmode
\small
\begin{array}[t]{lrll@{\hskip1ex}l}
\desc{Program} &\Pi      & ::= &  \multicolumn{2}{l}{\rep{\psi}, e}\\
\desc{Class declaration} &\psi     & ::= &  \multicolumn{2}{l}{\traitdecl{T}{\rep{V\;\beta}}{\rep{t}}{\rep{\mu}}                     
~\mid~\objectdecl{O}{\rep{\beta}}{\rep{x\COLON \tau}}{\rep{t}}{\rep{\mu}}}\\
\desc{Method definition} &\mu      & ::= &  \multicolumn{2}{l}{\methoddef{m}{\rep{\kappa}}{\rep{x\COLON\tau}}{\tau}{e}}\\
\desc{Class type parameter binding} &\beta    & ::= &  \multicolumn{2}{l}{\bindup{P}{\rep{\tau}}}\\
\desc{Method type parameter binding} &\kappa   & ::= &  \multicolumn{2}{l}{\lobindup{\rep{\tau}}{P}{\rep{\tau}}}\\
\desc{Variance mark} &V        & ::= &  \bindco &~\mid~ \bindcontra~\mid~\bindinv\\
\desc{Expression} &e      & ::= &  z
&~\mid~ ((\bar{x\COLON\tau})\COLON\tau \Rightarrow e)
~\mid~ e\apply(\bar{e})
~\mid~  O\obb{\tau}(\bar{e})
~\mid~  e.m(\rep{e})\\
\desc{Bindable variable} &z        & ::= &  x &~\mid~ \kwd{self}\\
\desc{Type} &\tau & ::= &  P
&~\mid~ c
~\mid~(\rep{\tau})
~\mid~\arrowtype{\tau}{\tau}
~\mid~\Any\\
\desc{Constructed type} &c        & ::= &  t
&~\mid~O\obb{\tau}\\
\desc{Trait type} &t        & ::= &  \multicolumn{2}{l}{T\obb{\tau}}
\end{array}
\vspace*{-.5em}
\caption{Abstract syntax of FGFV}
\label{fig:grammar}
\end{figure}

Figure~\ref{fig:grammar} shows the syntax of the FGFV calculus.
Metavariables $T$ and $U$ range over trait names;
$O$ ranges over object names;
$x$ ranges over variable and field names;
$z$ is either $x$ or the keyword $\KWD{self}$;
$m$ ranges over method names;
$P$, $Q$, $S$ range over type parameter names;
and $\tau$, $\omega$, $\zeta$, $\xi$ range over types.

A program consists of a sequence of class declarations and
a top-level expression.
A class is either a trait or an object;
traits are like Java interfaces that do not have any fields, and
objects are like Java final classes that cannot be extended.
A trait may have type parameters with variance annotations and
upper bounds, extend any number of other traits, and have method definitions.
An object is similar to a trait except that (1)~its type parameters do
not have variance annotations because they are always invariant, and
(2)~it may have fields, which are initialized when an object instance is
created.  When a class does not extend any trait, we omit the
empty extends clause $\extends \bigbd{}$.
When a class has no type parameters, we omit the
empty brackets; we write $T$ for $T\ob{}$.
A method may have type parameters as well, but they have both
lower and upper bounds.
We omit clauses when no bounds are specified; for example,
we write $T\ob{\bindco \; P}$ for $T\ob{\bindco \; \bindup{P}{}}$ and
$m\ob{Q,\bd{Q} \extends S}$ for $m\ob{\lobindup{}{Q}{}, \lobindup{Q}{S}{}}$.
By default, each type parameter has an upper bound $\Any$ and a
lower bound $\Bottom$, which we describe in the next subsection.

The bounds of a type parameter place restrictions on valid instantiations
of the type parameter; a valid type instantiation must be a subtype of
all its upper bounds and, at the same time, a supertype of
all its lower bounds, if any.
We write $\bindco$ for covariance, $\bindcontra$ for
contravariance, and $\bindinv$ for invariance.  As discussed in
Section~\ref{sec:motiv:fgfv}, variance annotations of type parameters define
subtype relations between different instances of a generic trait.
For example, consider the following trait declarations:\\[-.5em]
\[
\begin{array}{l}
\kwd{trait} \; U_1 \; \dontcares \; \kwd{end}\\
\kwd{trait} \; U_2 \extends \bigbd{U_1} \; \dontcares \; \kwd{end}\\
\kwd{trait} \; T \bigob{\bindco \; P, \bindcontra \; Q, \bindinv \; S} \;
\dontcares \; \kwd{end}
\end{array}
\]
where ``$\dontcares$'' denotes omitted material.
Because $U_2$ is a subtype of $U_1$,
the variance annotations of the type parameters of $T$ define
subtype relations such as $T\ob{U_2, U_1, U_1} \subtypeof T\ob{U_1, U_2, U_1}$.

\begin{figure}[t]
\small
\newjudge{Well-formed programs}{\jtype[]{\Pi}{g}}

\infrule[T-Program]
  { \Pi = \bar{\psi},e  \andalso
    \mathit{distinct}(\rep{\mathit{name}(\psi)})  \andalso
    \Delta = \setb{\psi}  \andalso
    \rep{\jwfdecl[\same{\Delta}]{\psi}}  \andalso
    \jtype[\Delta;\emptyseq]{e}{\expty{\dontcare}{g}} }
  { \jtype[]{\Pi}{g} }

\newjudge{Well-formed trait and object declarations}{\jwfdecl{\psi}}

{\showboxdepth=20\showboxbreadth=20

\infrule[D-Trait]
  { \Delta' = \Delta \cup \bigsetb{\lobindup{}{P}{\rep{\xi}}}  \andalso
    \mathit{distinct}(\rep{P})  \andalso
    \rep{\rep{\jwftype[\same{\same{\Delta'}}]{\xi}}}  \andalso
    \rep{\jwftype[\same{\Delta'}]{t}}  \\[2.5pt]
    \setb{J} = \properancestors[\Delta']{T\obb{P}}  \andalso
    \rep{\jancestormeetrule[\Delta'; \rep{J}]{J_i}{J_j}}[1 \leq i < j \leq \countofb{J}] \\[3.5pt]
    \rep{\same{T} \neq \mathit{name}(J)}  \andalso
    \rep{\jwfmeth[\same{\Delta';\kwd{self}\COLON T\bigobb{P};\rep{V\;P}}]{\mu}}  \andalso
    \setb{d} = \allvisible[\Delta']{T\obb{P}}  \\[3.5pt]
    \rep{\jnoduplicatesrule[\Delta']{d_i}{d_j}}[1 \leq i < j \leq \countofb{d}]  \andalso
    \rep{\jmethodmeetrule[\Delta']{d_i}{d_j}{\setb{d}}}[1 \leq i < j \leq \countofb{d}] \\[3pt]
    \rep{\jreturntyperule[\Delta']{d_i}{d_j}}[1 \leq i < j \leq \countofb{d}] \andalso
    \rep{\jreturntyperule[\Delta']{d_j}{d_i}}[1 \leq i < j \leq \countofb{d}]
    }
  { \jwfdecl{\traitdecl{T}{\rep{\bindup{V\;P}{\rep{\xi}}}}{\rep{t}}{\rep{\mu}}} }

}

\infrule[D-Object]
  { \Delta' = \Delta \cup \bigsetb{\lobindup{}{P}{\rep{\xi}}}  \andalso
    \mathit{distinct}(\bar{P})  \andalso
    \rep{\rep{\jwftype[\same{\same{\Delta'}}]{\xi}}}  \andalso
    \rep{\jwftype[\same{\Delta'}]{t}}  \\[2.5pt]
    \setb{J} = \properancestors[\Delta']{O\obb{P}}  \andalso
    \rep{\jancestormeetrule[\Delta'; \rep{J}]{J_i}{J_j}}[1 \leq i < j \leq \countofb{J}]  \\[3.5pt]
    \rep{\jwfmeth[\same{\Delta';\kwd{self}\COLON O\bigobb{P},\bar{x\COLON\tau};\rep{\bindinv\;P}}]{\mu}} \andalso
    \setb{d} = \allvisible[\Delta']{O\obb{P}}  \\[3.5pt]
    \rep{\jnoduplicatesrule[\Delta']{d_i}{d_j}}[1 \leq i < j \leq \countofb{d}]  \andalso
    \rep{\jmethodmeetrule[\Delta']{d_i}{d_j}{\setb{d}}}[1 \leq i < j \leq \countofb{d}] \\[3pt]
    \rep{\jreturntyperule[\Delta']{d_i}{d_j}}[1 \leq i < j \leq \countofb{d}] \andalso
    \rep{\jreturntyperule[\Delta']{d_j}{d_i}}[1 \leq i < j \leq \countofb{d}] \\[3pt]
    \jbwftype[\Delta']{\tau}  \andalso
    \mathit{distinct}(\bar{x})
    }
  { \jwfdecl{\objectdecl{O}{\rep{\bindup{P}{\rep{\xi}}}}{\rep{x\COLON\tau}}{\rep{t}}{\rep{\mu}}} }

\newjudge{Well-formed method declarations}{\jwfmeth{\mu}}

\infrule[D-Method]
  {
    \jscope{\bigobb{\lobindup{\rep{\zeta}}{Q}{\rep{\xi}}}} \andalso
    \Delta' = \Delta \cup \bigsetb{\lobindup{\rep{\zeta}}{Q}{\rep{\xi}}} \andalso
\jbbwftype[\Delta']{\zeta} \andalso
    \jbbwftype[\Delta']{\xi} \\[3pt]
    \mathit{distinct}(\bar{x})  \andalso
    \rep{\jwftype[\same{\Delta'}]{\tau}}  \andalso  \jwftype[\Delta']{\omega} \andalso
    \jtype[\Delta';\Gamma,\bar{x\COLON\tau}]{e}{\expty{\dontcare}{\rho}}  \andalso
\jsubtype[\Delta']{\rho}{\omega} \\[3pt]
    \mathit{distinct}(\bar{P},\bar{Q})  \andalso
    \rep{(V = \bindinv)\lor(\forall \T.\ (\same{\arrowtype{(\rep{\tau})}{\omega}}\neq\T\context{P}) \lor (\jcontextvar[\same{\Delta}]{\T}{V}))} }
  { \jwfmeth{\methoddef{m}{\rep{\lobindup{\rep{\zeta}}{Q}{\rep{\xi}}}}{\rep{x\COLON\tau}}{\omega}{e}} }

\vspace*{-1.5em}
\caption{Well-formed programs, class declarations, and method declarations}
\label{fig:wellformed}
\end{figure}

Expressions and types are conventional.
One can refer to a field or bound variable simply by mentioning it;
the keyword $\KWD{self}$ refers to the target object for the currently executing method.
Functions are first-class values; $((\bar{x\COLON\tau})\COLON\tau \Rightarrow e)$
is a (typed) function constructor,
and we write $e\apply(\bar{e})$ for
function calls.  An object creation $O\obb{\tau}(\bar{e})$ makes an
instance of $O$ by instantiating its type parameters with $\bar{\tau}$ and
initializing its fields with the values of $\bar{e}$.
While tuples are not first-class values in FGFV (they would be trivial to add,
but are not needed for our purposes), we use tuple types $(\rep{\tau})$
to denote the domain types of functions and methods, which may have multiple parameters.
The type $\Any$ is a supertype of every type (and $\Bottom$, used below, is a subtype of every type but is intentionally not expressible in the surface syntax).


\subsection{Well-formedness}\label{sec:wellformed}

We define well-formedness of a program and its declarations as shown in
Figure~\ref{fig:wellformed}.
A type-declaration environment $\Delta$ is a set of class declarations and
type parameter bindings, and
a variable-type environment $\Gamma$ is a sequence of bindable variables
and their types.  Checking well-formedness may introduce
\textit{internal types}:
\begin{center}
$\descsm{Internal type} \quad \alpha~::= P~\mid~C~\mid~\tupleb{\alpha}~\mid~
\arrowtype{\alpha}{\alpha}~\mid~\Any~\mid~\Bottom~\mid~
\uniontype{\alpha}{\alpha}~\mid~\intersectiontype{\alpha}{\alpha}$
\end{center}
which contain $\Bottom$, union types, and intersection types.
Type $\Bottom$ is a subtype of every type; hence, $\Any$ and $\Bottom$
are the top and bottom, respectively, of the type lattice.
Metavariables $\alpha, \gamma, \rho, \chi, \eta$ range over internal types.
We introduce metavariables $K$, $C$, $J$, and $M$,
analogous to $\kappa$, $c$, $t$, and $\mu$, respectively, where 
types $\tau$ in the latter are replaced by internal types $\alpha$ in the former.

Rule \rncite{T-Program} states that a program is well-formed
if its class declarations have distinct names and are well-formed, and
its top-level expression is well-formed.
%%% The following material has been obsoleted by new section on notation.
% We write $\emptyseq$ for an empty sequence.
% Following a suggestion of Steele~\citeyearpar[0:59:37]{NEW-OLD-LANGUAGE},
% we use
% an underline below submaterial within overlined material to prevent the 
% overline from attaching subscripts to the metavariables in the submaterial;
% hence, $\rep{\jwfdecl[\same{\Delta}]{\psi}}$ represents the sequence
% $\Delta\vdash\psi_1\, \mathsf{ok}\ \ \Delta\vdash\psi_2\, \mathsf{ok}\ \cdots\ \Delta\vdash\psi_n\, \mathsf{ok}$
% for any appropriate $n \geq 0$.
The type of a program is the type of its top-level expression;
the metavariable $g$ and $h$ represents a \textit{ground type},
which is a type that contains no type parameters.
We describe well-formedness and type checking of expressions
in Section~\ref{sec:static}.

Rules \rncite{D-Trait} and \rncite{D-Object} check well-formedness
of class declarations.  For the declaration of a class $C$,
they verify that the type parameter names of the declaration are distinct,
that the types used in the declaration are well-formed,
and that the declared methods are well-formed.
In addition, as we explain later in this subsection,
they make sure that (1)~the \textit{proper ancestors} of $C$,
which are all the traits $C$ extends, by explicitly extending and
transitively inheriting as well, are well-formed, and
(2)~the \textit{visible methods} of $C$,
which are the methods either declared by $C$ or inherited
from the traits $C$ extends, obey the overloading rules.
The only differences are that rule \rncite{D-Trait} additionally
guarantees that a trait does not extend itself, and rule \rncite{D-Object}
guarantees that the names of the fields are distinct.

Well-formedness of types is conventional
as shown in Figure~\ref{fig:wellformedtypes}.
The types $\Any$ and $\Bottom$ are trivially well-formed.
A type parameter $P$ is well-formed under a type-declaration environment
$\Delta$, if $\Delta$ contains a binding of $P$.
A constructed type $T\obb{\alpha}$ or $O\obb{\alpha}$ is well-formed, if
the trait $T$ or the object $O$ is declared in the given
type-declaration environment and all the type arguments $\rep{\alpha}$
are well-formed and satisfy the declared
bounds of their corresponding type parameters.
Finally, the remaining structural types such as tuple types, arrow types, 
union types, and intersection types are well-formed,
if all their component types are.

\begin{figure}[t]
\small
\newjudge{Well-formed types}{\jwftype{\alpha}}
\smallskip
\begin{tabular}{@{}c@{\hskip0.02\textwidth}c@{}}
\begin{minipage}[t]{0.48\textwidth}
\infax[W-Any]
  { \jwftype{\Any} }
\smallskip

\infax[W-Bottom]
  { \jwftype{\Bottom} }
\smallskip

\infrule[W-Param]
  { \lobindup{\dontcares}{P}{\dontcares} \in \Delta }
  { \jwftype{P} }
\smallskip

\infrule[W-Tuple]
  { \jbwftype{\alpha} }
  { \jwftype{(\bar{\alpha})} }

\end{minipage}&
\begin{minipage}[t]{.5\textwidth}
\infrule[W-Arrow]
  { \jwftype{\alpha}  \andalso  \jwftype{\rho} }
  { \jwftype{\arrowtype{\alpha}{\rho}} }
\smallskip

\infrule[W-Union]
  { \jwftype{\alpha}  \andalso  \jwftype{\gamma} }
  { \jwftype{\uniontype{\alpha}{\gamma}} }
\smallskip

\infrule[W-Intersection]
  { \jwftype{\alpha}  \andalso  \jwftype{\gamma} }
  { \jwftype{\intersectiontype{\alpha}{\gamma}} }
\smallskip

\end{minipage}
\end{tabular}
\infrule[W-Trait]
  { \traitdecl{T}{\rep{\bindup{V\;P}{\rep{\xi}}}}{\dontcares}{\dontcares} \in \Delta \andalso
    \jbwftype{\alpha} \andalso
    \jbsubtype{\alpha}{\same{\substb{\alpha}{P}}\bigtypeintersection\setb{\xi}} }
  { \jwftype{T\obb{\alpha}} }
\smallskip

\infrule[W-Object]
  { \objectdecl{O}{\rep{\bindup{P}{\rep{\xi}}}}{\dontcares}{\dontcares}{\dontcares} \in \Delta \andalso
    \jbwftype{\alpha} \andalso
    \jbsubtype{\alpha}{\same{\substb{\alpha}{P}}\bigtypeintersection\setb{\xi}} }
  { \jwftype{O\obb{\alpha}} }
\vspace*{-1em}
\caption{Well-formed types}
\label{fig:wellformedtypes}
\end{figure}

The proper ancestors $\setb{J}$ of a class $C$ are well-formed, if
every pair in $\setb{J}$ satisfies one of the well-formed ancestors rules
shown in Figure~\ref{fig:ancestors}: $\rep{\jancestormeetrule
[\Delta'; \rep{J}]{J_i}{J_j}}[1 \leq i < j \leq \countofb{J}]$.
The well-formed ancestors rules ensure that if a class $C$ extends
two instances of a trait $T$, $T\obb{\alpha}$ and $T\obb{\rho}$,
then $C$ also extends some $T\obb{\gamma}$
(which may be $T\obb{\alpha}$, $T\obb{\rho}$,
or a third instance) that extends both $T\obb{\alpha}$ and $T\obb{\rho}$.
Therefore, because the rules guarantee that there always exists a single
minimal ancestor when a class extends multiple instances of a trait,
no ambiguous method calls are possible due to multiple instances of a trait
as discussed in Section~\ref{sec:dynamic}.


\begin{figure}[t]
\small
\newjudge{Well-formed ancestors}{\jancestormeetrule{J}{J}}
\vspace*{-.7em}

\infrule[Anc-Diff-Trait]
  { T \neq T' }
  { \jancestormeetrule{T\obb{\alpha}}{T'\obb{\eta}} }

\vspace*{-1em}
\infrule[Anc-Same-Trait]
  { T\obb{\gamma} \in \setb{J}  \andalso
    \jsubtype{T\obb{\gamma}}{T\obb{\alpha}}  \andalso
    \jsubtype{T\obb{\gamma}}{T\obb{\eta}} }
  { \jancestormeetrule{T\obb{\alpha}}{T\obb{\eta}} }

\newjudge{Well-formed type parameter bindings}{\jscope{\obb{K}}}
\vspace*{-.7em}

\infax[Binding-Empty]
  { \jscope{\ob{}} }

\vspace*{-1.5em}
\infrule[Binding-Step]
 { \FV(\bar{\chi'})\subseteq\mathit{parameters}(\Delta)\cup\setb{P} \andalso
   \FV(\bar{\eta'})\subseteq\mathit{parameters}(\Delta) \\%[3pt]
   \jsubtype{\bigtypeunion\bigsetb{\chi'} }{ \bigtypeintersection\bigsetb{\eta'} } \andalso
   \jscope{\obb{\lobindup{\rep{\chi}}{P}{\rep{\eta}}}} }
 { \jscope{\ob{\rep{\lobindup{\rep{\chi}}{P}{\rep{\eta}}},\lobindup{\rep{\chi'}}{P'}{\rep{\eta'}}}} }

\vspace*{-1.5em}
\caption{Well-formed ancestors and type parameter bindings}
\label{fig:ancestors}
\end{figure}

The visible methods $\setb{d} = \setb{\tuple{C,M}}$ of a class 
obey the overloading rules if every pair in $\setb{d}$ satisfies 
the overloading rules described in Section~\ref{sec:rules:formal}.
The set $\setb{\tuple{C,M}}$ contains pairs of visible method declarations 
$\bar{M}$ and their defining classes $\bar{C}$.

Rule \rncite{D-Method} checks well-formedness of method declarations.
For the method declaration $\mu$, it verifies that the parameter names
of the declaration are distinct;
that the class type parameters and the method type parameters
of the declaration are collectively distinct;
that types used in the declaration are well-formed;
that the body expression is well-formed;
and that the type of the body expression is a subtype of the 
declared return type.  In addition, as we discuss below,
it checks that method type parameter bindings are well-formed and
that any references to trait type parameters are correctly used
according to their variance annotations.

A sequence of method type parameter bindings is well-formed,
if (1)~every reference to a method type parameter falls within its scope
(which extends to the right of its declaration),
(2)~no reference to a method type parameter appears in any of the upper bounds,
and (3) for each type parameter, the union of its lower bounds is a subtype
of the intersection of its upper bounds.
As we discuss in Section~\ref{sec:discuss},
restriction (1) makes dynamic dispatch more efficient
and restriction (2) makes it consistent.  Figure~\ref{fig:ancestors}
shows the rules for well-formed type parameter bindings.

Trait type parameters are correctly used according to their variance annotations,
if covariant and contravariant type parameters appear only at
covariant and contravariant positions, respectively.
In order to express this restriction, we introduce a \textit{type context} $\T$,
which defines positions where a type can appear within another type.
Similarly to the conventional evaluation contexts~\cite{wright94},
a type context $\T$ is a type with one of its components replaced by a hole,
denoted by ``$\emptycontext$''.  The type $\T\context{\tau}$ results from
placing a type $\tau$ in the hole of $\T$:
\begin{center}
$\descsm{Type context} \quad 
\T~::=~\emptycontext ~\mid~ T\ob{\bar\tau, \T, \bar\tau} ~\mid~ O\ob{\bar\tau, \T, \bar\tau} ~\mid~ \tuple{\bar\tau, \T, \bar\tau} ~\mid~ \arrowtype{\T}{\tau} ~\mid~\arrowtype{\tau}{\T}$
\end{center}
Thus, if $\tau = \arrowtype{\tuple{P,T\ob{Q}}}{\Any}$, 
then $\tau = \T\context{P} = \T'\context{Q}$ where 
$\T = \arrowtype{\tuple{\emptycontext, T\ob{Q}}}{\Any}$ and 
$\T' = \arrowtype{\tuple{P, T\ob{\emptycontext}}}{\Any}$.

\begin{figure}[t]
\small
\smallskip
\newjudge{Variance of type context}{\jcontextvar{\T}{V}}
\begin{tabular}{@{}c@{\hskip0.01\textwidth}c@{}}
\begin{minipage}[t]{0.5\textwidth}
\infax[Tc-Empty] {\jcontextvar{\emptycontext}{\bindco}}
\smallskip

\infrule[Tc-Trait]
  { \traitdecl{T}{\rep{\bindup{V\;\dontcare}{\dontcares}}}{\dontcares}{\dontcares} \in \Delta \\[3pt]
    i = \countof{\rep\tau} + 1  \andalso
    \countofb{V} = \countof{\rep\tau} + \countof{\rep\omega} + 1  \\[3pt]
    \jcontextvar{\T}{V'} \andalso
    V = V_i \times V' }
  { \jcontextvar{T\ob{\bar\tau,\T,\bar\omega}}{V} }
\smallskip

\infax[Tc-Obj] { \jcontextvar{O\ob{\bar\tau,\T,\bar\omega}}{\bindinv} }

\smallskip

\newjudge{Flip of variance}{\flip{V} = V}
\begin{equation*}
\flip{V} =
\begin{cases}
\bindcontra & \text{if } V = \bindco\\
\bindco & \text{if } V = \bindcontra\\
\bindinv & \text{if } V = \bindinv\\
\end{cases}
\end{equation*}

\end{minipage}&
\begin{minipage}[t]{.5\textwidth}
\infrule[Tc-Tuple]
  { \jcontextvar{\T}{V} }
  { \jcontextvar{(\bar\tau,\T,\bar\omega)}{V} }
\smallskip

\infrule[Tc-Arrow-Param]
  { \jcontextvar{\T}{V'}
  { \andalso V=\flip{V'} } }
  { \jcontextvar{\arrowtype{\T}{\tau}}{V} }
\smallskip

\infrule[Tc-Arrow-Return]
  { \jcontextvar{\T}{V} }
  { \jcontextvar{\arrowtype{\tau}{\T}}{V} }
\smallskip

\smallskip

\newjudge{Composition of variance}{V \times V = V}

\begin{equation*}
V_1 \times V_2 =
\begin{cases}
\bindinv & \text{if } V_1 = \bindinv \lor V_2 = \bindinv\\
\bindco & \text{if } V_1 \neq \bindinv \land V_2 \neq \bindinv \land V_1 = V_2\\
\bindcontra & \text{if } V_1 \neq \bindinv \land V_2 \neq \bindinv \land V_1 \neq V_2\\
\end{cases}
\end{equation*}

\end{minipage}
\end{tabular}
\vspace*{-1em}
\caption{Variance of type context}
\label{fig:contextvar}
\end{figure}

Using type contexts, the last premise of rule \rncite{D-Method} ensures that
the class type parameters of a method appear in correct positions.
Considering the arrow type representing the method's type
$\tau' = \arrowtype{(\rep{\tau})}{\omega}$,
an invariant type parameter can appear anywhere in $\tau'$. However, 
when a covariant or contravariant type parameter $P$ appears in 
$\tau' = \T\context{P}$, the \textit{variance of the type context} of $P$
should match the variance $V$ of $P$: $\Delta\vdash{\T}\,\mathsf{variance}\,V$.
As shown in Figure~\ref{fig:contextvar},
the variance of a type context is defined inductively:
(1)~the hole is covariant;
(2)~components of tuple types, return types of arrow types, and
covariant type parameters of traits preserve
the variance of their enclosing contexts;
(3)~parameter types of arrow types and
contravariant type parameters of traits flip
the variance of their enclosing contexts; and
(4)~type parameters of objects and
invariant type parameters of traits are invariant.
For example, if $\tau = \arrowtype{\tuple{P,T\ob{Q}}}{\Any}$
is a method's type where $P$ and $Q$ are class type parameters and
$T$ has a contravariant type parameter,
then
$P$ is required to be contravariant or invariant and $Q$ is required to be covariant or invariant.
